% @file paper44_gruppentheorie_de.tex
% @brief Paper 44 (DE): Gruppentheorie — Endliche Gruppen, Sylow-Theorie,
%        Klassifikation und Anwendungen
% @author Michael Fuhrmann
% @date 2026-03-12
% @build 135

\documentclass[12pt,a4paper]{amsart}
\usepackage{amsmath,amssymb,amsthm}
\usepackage{mathtools}
\usepackage[utf8]{inputenc}
\usepackage[T1]{fontenc}
\usepackage[ngerman]{babel}
\usepackage{hyperref}
\usepackage{enumitem,booktabs,array}
\usepackage{geometry}
\geometry{margin=2.5cm}

%% Theorem-Umgebungen auf Deutsch
\newtheorem{theorem}{Satz}[section]
\newtheorem{lemma}[theorem]{Lemma}
\newtheorem{corollary}[theorem]{Korollar}
\newtheorem{proposition}[theorem]{Proposition}
\newtheorem{conjecture}[theorem]{Vermutung}
\theoremstyle{definition}
\newtheorem{definition}[theorem]{Definition}
\newtheorem{example}[theorem]{Beispiel}
\newtheorem{remark}[theorem]{Bemerkung}

\title{Gruppentheorie: Endliche Gruppen, Sylow-Theorie,\\
       Klassifikationss\"atze und Anwendungen}

\author{Michael Fuhrmann}
\address{Specialist Mathematics Project, Build 135}
\email{specialist-maths@localhost}
\date{\today}

\begin{document}

\begin{abstract}
Die Gruppentheorie ist die mathematische Theorie der Symmetrie und ihrer algebraischen
Abstraktion. Dieses Paper gibt eine umfassende Darstellung der endlichen Gruppentheorie.
Ausgehend von den grundlegenden Definitionen und den drei Isomorphiesätzen werden zyklische
und abelsche Gruppen untersucht, bis hin zum vollständigen Klassifikationssatz für endliche
abelsche Gruppen. Gruppenoperationen, der Bahn-Stabilisator-Satz und das Burnside-Lemma
werden vollständig entwickelt, gefolgt von den Sylow-Sätzen mit Beweisen. Auflösbare und
nilpotente Gruppen werden im Zusammenhang mit Kompositionsreihen und dem Satz von
Jordan--Hölder behandelt, mit Anwendungen auf die Unlösbarkeit des Quintikums
(Abel--Ruffini). Die Theorie einfacher Gruppen wird diskutiert, einschließlich des Satzes
von Feit--Thompson und der Klassifikation endlicher einfacher Gruppen (CFSG). Der
abschließende Abschnitt gibt einen Überblick über Anwendungen in der Galois-Theorie,
Kristallographie und Quantenmechanik.
\end{abstract}

\maketitle
\tableofcontents

%% ============================================================
\section{Einleitung}
\label{sec:einleitung}
%% ============================================================

Der Begriff der \emph{Gruppe} entstand im frühen neunzehnten Jahrhundert durch die Arbeit
von \'{E}variste Galois, der das einführte, was wir heute Permutationsgruppen nennen, um zu
charakterisieren, welche Polynomgleichungen durch Radikale lösbar sind. Unabhängig davon
definierte Arthur Cayley 1854 abstrakte Gruppen und erkannte, dass dieselbe algebraische
Struktur Permutationen, ganzen Zahlen unter Addition, Symmetrien geometrischer Objekte und
einer Vielzahl weiterer mathematischer Phänomene zugrunde liegt.

Heute ist die Gruppentheorie ein Grundpfeiler der reinen Mathematik und der theoretischen
Physik. Ihre Anwendungen umfassen:

\begin{itemize}[leftmargin=2em]
  \item \textbf{Symmetrieanalyse}: Jede physikalische Symmetrie entspricht einer Gruppe;
        Erhaltungsgrößen entstehen gemäß dem Noether-Theorem aus Symmetriegruppen der
        Wirkung.
  \item \textbf{Galois-Theorie}: Die Lösbarkeit einer Polynomgleichung durch Radikale ist
        äquivalent zur Auflösbarkeit ihrer Galois-Gruppe.
  \item \textbf{Kryptographie}: Das diskrete Logarithmusproblem in endlichen zyklischen
        Gruppen und elliptischen Kurvengruppen bildet die Grundlage moderner
        Public-Key-Kryptographie.
  \item \textbf{Kristallographie}: Die 17 Wandgruppen und 230 Raumgruppen klassifizieren
        alle möglichen Kristallsymmetrien in der Ebene und im dreidimensionalen Raum.
  \item \textbf{Quantenmechanik}: Die Darstellungstheorie von Lie-Gruppen wie
        $\mathrm{SU}(2)$ und $\mathrm{SO}(3)$ bestimmt die Struktur atomarer Spektren
        und den Spin von Elementarteilchen.
\end{itemize}

Das Paper ist wie folgt gegliedert. Abschnitt~\ref{sec:definitionen} legt die grundlegenden
Definitionen fest. Abschnitt~\ref{sec:isomorphie} präsentiert die drei Isomorphiesätze.
Die Abschnitte~\ref{sec:zyklisch} und~\ref{sec:operationen} behandeln zyklische und abelsche
Gruppen bzw.\ Gruppenoperationen mit Sylow-Theorie. Die
Abschnitte~\ref{sec:aufloesbar} und~\ref{sec:einfach} befassen sich mit auflösbaren und
nilpotenten Gruppen sowie einfachen Gruppen. Abschnitt~\ref{sec:anwendungen} gibt einen
Überblick über die wichtigsten Anwendungen. Abschnitt~\ref{sec:schluss} schließt ab.

%% ============================================================
\section{Grunddefinitionen}
\label{sec:definitionen}
%% ============================================================

\begin{definition}[Gruppe]
Eine \emph{Gruppe} ist ein Paar $(G, \cdot)$, bestehend aus einer nichtleeren Menge $G$
und einer binären Verknüpfung $\cdot : G \times G \to G$, die folgende Axiome erfüllt:
\begin{enumerate}[label=(\roman*)]
  \item \textbf{Assoziativität}: $(a \cdot b) \cdot c = a \cdot (b \cdot c)$ für alle
        $a,b,c \in G$.
  \item \textbf{Neutrales Element}: Es gibt ein $e \in G$ mit $e \cdot a = a \cdot e = a$
        für alle $a \in G$.
  \item \textbf{Inverse}: Für jedes $a \in G$ existiert ein $a^{-1} \in G$ mit
        $a \cdot a^{-1} = a^{-1} \cdot a = e$.
\end{enumerate}
Gilt zusätzlich $a \cdot b = b \cdot a$ für alle $a, b \in G$, so heißt die Gruppe
\emph{abelsch} (oder \emph{kommutativ}). Die \emph{Ordnung} $|G|$ ist die Mächtigkeit
von~$G$.
\end{definition}

\begin{example}
Die ganzen Zahlen $(\mathbb{Z}, +)$ bilden eine unendliche abelsche Gruppe; der
Einheitskreis $\mathrm{U}(1) = \{z \in \mathbb{C} : |z| = 1\}$ unter Multiplikation ist
eine unendliche abelsche Gruppe; die symmetrische Gruppe $S_n$ aller Permutationen von
$\{1, \ldots, n\}$ unter Komposition ist eine endliche Gruppe der Ordnung~$n!$.
\end{example}

\begin{definition}[Untergruppe]
Eine nichtleere Teilmenge $H \subseteq G$ heißt \emph{Untergruppe} von $G$, geschrieben
$H \leq G$, wenn $H$ selbst eine Gruppe unter der eingeschränkten Verknüpfung von $G$ ist.
Äquivalent gilt $H \leq G$ genau dann, wenn für alle $a, b \in H$ gilt: $ab^{-1} \in H$.
\end{definition}

\begin{definition}[Nebenklasse und Index]
Für $H \leq G$ und $a \in G$ ist die \emph{linke Nebenklasse} von $H$ bezüglich $a$
die Menge $aH = \{ah : h \in H\}$. Die Anzahl der verschiedenen linken Nebenklassen
heißt \emph{Index} $[G : H]$.
\end{definition}

\begin{theorem}[Lagrange, 1771]
\label{thm:lagrange}
Sei $G$ eine endliche Gruppe und $H \leq G$. Dann teilt $|H|$ die Ordnung $|G|$, und es
gilt:
\[
  |G| = |H| \cdot [G : H].
\]
\end{theorem}

\begin{proof}
Die linken Nebenklassen von $H$ zerlegen $G$ in $[G:H]$ disjunkte Teilmengen, jede
von Mächtigkeit $|H|$. Abzählen ergibt $|G| = [G:H] \cdot |H|$.
\end{proof}

\begin{corollary}
Die Ordnung jedes Elements $a \in G$ teilt $|G|$.
\end{corollary}

\begin{definition}[Normalteiler und Quotientengruppe]
Eine Untergruppe $N \leq G$ heißt \emph{Normalteiler}, geschrieben $N \trianglelefteq G$,
wenn $gNg^{-1} = N$ für alle $g \in G$. Ist $N \trianglelefteq G$, so bildet die Menge
der Nebenklassen $G/N = \{gN : g \in G\}$ unter der Verknüpfung $(aN)(bN) = (ab)N$
eine Gruppe, die \emph{Quotientengruppe} (oder \emph{Faktorgruppe}) $G$ nach $N$.
\end{definition}

\begin{definition}[Homomorphismus und Isomorphismus]
Eine Abbildung $\varphi : G \to H$ zwischen Gruppen heißt \emph{Homomorphismus}, wenn
$\varphi(ab) = \varphi(a)\varphi(b)$ für alle $a,b \in G$. Der \emph{Kern} von $\varphi$
ist $\ker\varphi = \{g \in G : \varphi(g) = e_H\}$ und das \emph{Bild} ist
$\mathrm{Im}\,\varphi = \{\varphi(g) : g \in G\}$.
Ein bijektiver Homomorphismus heißt \emph{Isomorphismus}; wir schreiben $G \cong H$,
wenn ein Isomorphismus zwischen $G$ und $H$ existiert.
\end{definition}

\begin{proposition}
Ist $\varphi : G \to H$ ein Homomorphismus, so gilt $\ker\varphi \trianglelefteq G$
und $\mathrm{Im}\,\varphi \leq H$.
\end{proposition}

%% ============================================================
\section{Die Isomorphiesätze}
\label{sec:isomorphie}
%% ============================================================

Die drei Isomorphiesätze gehören zu den nützlichsten Strukturresultaten der Gruppentheorie.
Sie wurden in ihrer modernen Form von Emmy Noether im Jahr 1927 formuliert.

\begin{theorem}[Erster Isomorphiesatz (Homomorphiesatz)]
\label{thm:erster_iso}
Sei $\varphi : G \to H$ ein Gruppenhomomorphismus. Dann gilt $\ker\varphi \trianglelefteq G$
und
\[
  G / \ker\varphi \;\cong\; \mathrm{Im}\,\varphi.
\]
Der Isomorphismus ist explizit gegeben durch $g\ker\varphi \mapsto \varphi(g)$.
\end{theorem}

\begin{proof}
Wir wissen bereits, dass $K := \ker\varphi \trianglelefteq G$. Definiere
$\tilde\varphi : G/K \to H$ durch $\tilde\varphi(gK) = \varphi(g)$.
\textit{Wohldefiniertheit}: Ist $gK = g'K$, dann $g'^{-1}g \in K$, also
$\varphi(g'^{-1}g) = e_H$, woraus $\varphi(g) = \varphi(g')$ folgt.
\textit{Homomorphismus}: $\tilde\varphi((gK)(g'K)) = \varphi(gg') = \varphi(g)\varphi(g')$.
\textit{Injektivität}: $\tilde\varphi(gK) = e_H$ impliziert $g \in K$, also
$gK = K = e_{G/K}$.
\textit{Surjektivität} auf $\mathrm{Im}\,\varphi$ folgt aus der Definition.
\end{proof}

\begin{remark}
Der erste Isomorphiesatz reduziert die Klassifikation homomorpher Bilder von $G$ auf die
Klassifikation der Normalteiler von $G$.
\end{remark}

\begin{theorem}[Zweiter Isomorphiesatz]
\label{thm:zweiter_iso}
Sei $H \leq G$ und $N \trianglelefteq G$. Dann ist $HN = \{hn : h \in H,\, n \in N\}$
eine Untergruppe von $G$, $H \cap N \trianglelefteq H$, und es gilt:
\[
  H / (H \cap N) \;\cong\; HN / N.
\]
\end{theorem}

\begin{proof}
Man verifiziert, dass $HN$ mittels der Normalteiler-Eigenschaft von $N$ eine Untergruppe
ist. Die natürliche Abbildung $\varphi : H \to HN/N$ definiert durch $\varphi(h) = hN$ ist
ein surjektiver Homomorphismus mit Kern $H \cap N$. Das Ergebnis folgt aus
Satz~\ref{thm:erster_iso}.
\end{proof}

\begin{theorem}[Dritter Isomorphiesatz]
\label{thm:dritter_iso}
Sei $N \trianglelefteq G$ und $K \trianglelefteq G$ mit $K \leq N$. Dann gilt
$N/K \trianglelefteq G/K$ und:
\[
  (G/K) / (N/K) \;\cong\; G/N.
\]
\end{theorem}

\begin{proof}
Die Abbildung $\varphi : G/K \to G/N$ definiert durch $\varphi(gK) = gN$ ist ein
wohldefinierter surjektiver Homomorphismus mit Kern $N/K$. Anwendung von
Satz~\ref{thm:erster_iso} liefert das Ergebnis.
\end{proof}

\begin{corollary}[Korrespondenzsatz]
Sei $N \trianglelefteq G$. Es gibt eine bijektive, inklusionserhaltende Entsprechung
zwischen den Untergruppen von $G/N$ und den Untergruppen von $G$, die $N$ enthalten.
Diese Entsprechung bewahrt die Normalteiler-Eigenschaft.
\end{corollary}

%% ============================================================
\section{Zyklische und abelsche Gruppen}
\label{sec:zyklisch}
%% ============================================================

\begin{definition}[Zyklische Gruppe]
Eine Gruppe $G$ heißt \emph{zyklisch}, wenn ein $g \in G$ existiert, sodass jedes Element
von $G$ eine Potenz von $g$ ist: $G = \langle g \rangle = \{g^k : k \in \mathbb{Z}\}$.
\end{definition}

\begin{theorem}
Jede zyklische Gruppe ist abelsch. Darüber hinaus gilt:
\begin{enumerate}[label=(\roman*)]
  \item Jede unendliche zyklische Gruppe ist isomorph zu $(\mathbb{Z}, +)$.
  \item Jede endliche zyklische Gruppe der Ordnung $n$ ist isomorph zu
        $\mathbb{Z}/n\mathbb{Z}$.
\end{enumerate}
\end{theorem}

\begin{theorem}
Jede Untergruppe einer zyklischen Gruppe ist zyklisch. Hat $G = \langle g \rangle$ die
Ordnung $n$, so sind die Untergruppen von $G$ genau die $\langle g^d \rangle$ für $d \mid n$,
und jede hat Ordnung $n/d$. Insbesondere besitzt $G$ für jeden Teiler $d$ von $n$ genau eine
Untergruppe der Ordnung $d$.
\end{theorem}

\begin{theorem}[Hauptsatz über endliche abelsche Gruppen]
\label{thm:hsuag}
Jede endliche abelsche Gruppe $G$ ist isomorph zu einem direkten Produkt zyklischer Gruppen
von Primzahlpotenzordnung:
\[
  G \;\cong\; \mathbb{Z}_{p_1^{a_1}} \times \mathbb{Z}_{p_2^{a_2}} \times \cdots
              \times \mathbb{Z}_{p_k^{a_k}},
\]
wobei $p_1^{a_1} p_2^{a_2} \cdots p_k^{a_k} = |G|$. Diese Zerlegung ist eindeutig bis auf
die Reihenfolge der Faktoren (Elementarteilerzerlegung).
\end{theorem}

\begin{example}
Die abelschen Gruppen der Ordnung $72 = 2^3 \cdot 3^2$ sind bis auf Isomorphie:
\begin{align*}
  &\mathbb{Z}_8 \times \mathbb{Z}_9,\quad
   \mathbb{Z}_4 \times \mathbb{Z}_2 \times \mathbb{Z}_9,\quad
   \mathbb{Z}_2^3 \times \mathbb{Z}_9,\\
  &\mathbb{Z}_8 \times \mathbb{Z}_3 \times \mathbb{Z}_3,\quad
   \mathbb{Z}_4 \times \mathbb{Z}_2 \times \mathbb{Z}_3 \times \mathbb{Z}_3,\quad
   \mathbb{Z}_2^3 \times \mathbb{Z}_3^2.
\end{align*}
Es gibt genau 6 nicht-isomorphe abelsche Gruppen der Ordnung 72.
\end{example}

\begin{theorem}[Chinesischer Restsatz für Gruppen]
Ist $\gcd(m,n) = 1$, so gilt $\mathbb{Z}_{mn} \cong \mathbb{Z}_m \times \mathbb{Z}_n$.
\end{theorem}

%% ============================================================
\section{Gruppenoperationen und Sylow-Theorie}
\label{sec:operationen}
%% ============================================================

\subsection{Gruppenoperationen}

\begin{definition}[Gruppenoperation]
Eine \emph{(linke) Operation} einer Gruppe $G$ auf einer Menge $X$ ist eine Abbildung
$G \times X \to X$, $(g,x) \mapsto g \cdot x$, die folgende Axiome erfüllt:
\begin{enumerate}[label=(\roman*)]
  \item $e \cdot x = x$ für alle $x \in X$,
  \item $(gh) \cdot x = g \cdot (h \cdot x)$ für alle $g,h \in G$, $x \in X$.
\end{enumerate}
\end{definition}

\begin{definition}[Bahn und Stabilisator]
Für $x \in X$ ist die \emph{Bahn} (oder \emph{Orbit}) von $x$ die Menge
$G \cdot x = \{g \cdot x : g \in G\}$, und der \emph{Stabilisator} (oder \emph{Isotropiegruppe})
von $x$ ist $G_x = \{g \in G : g \cdot x = x\} \leq G$.
\end{definition}

\begin{theorem}[Bahn-Stabilisator-Satz]
\label{thm:bahn_stab}
Sei $G$ eine endliche Gruppe, die auf $X$ operiert, und sei $x \in X$. Dann gilt:
\[
  |G \cdot x| \cdot |G_x| = |G|.
\]
\end{theorem}

\begin{theorem}[Burnside-Lemma, 1897]
\label{thm:burnside}
Sei $G$ eine endliche Gruppe, die auf einer endlichen Menge $X$ operiert. Die Anzahl
der verschiedenen Bahnen ist:
\[
  |X/G| = \frac{1}{|G|} \sum_{g \in G} |X^g|,
\]
wobei $X^g = \{x \in X : g \cdot x = x\}$ die Menge der Fixpunkte von $g$ bezeichnet.
\end{theorem}

\begin{example}[Halskettenzählung]
Wir zählen die Anzahl verschiedener Halsketten mit 6 Perlen, von denen jede in einer von
$c$ Farben gefärbt ist, wobei Drehungen als äquivalent betrachtet werden. Die
Drehgruppe $G = \mathbb{Z}_6$ operiert auf der Menge $X$ aller gefärbten Folgen
($|X| = c^6$). Nach dem Burnside-Lemma ist die Anzahl verschiedener Halsketten:
\[
  N(c) = \frac{1}{6}\bigl(c^6 + c^3 + 2c^2 + 2c\bigr).
\]
Für $c = 2$ ergibt sich $N(2) = \tfrac{1}{6}(64 + 8 + 8 + 4) = 14$.
\end{example}

\subsection{Sylow-Theorie}

Die Sylow-Sätze, bewiesen von Ludwig Sylow im Jahr 1872, sind das wichtigste Werkzeug
zum Verständnis der Untergruppenstruktur endlicher Gruppen.

\begin{definition}[Sylow-$p$-Untergruppe]
Sei $G$ eine endliche Gruppe und $p$ eine Primzahl mit $p^k \mid |G|$, aber
$p^{k+1} \nmid |G|$. Eine Untergruppe $P \leq G$ mit $|P| = p^k$ heißt
\emph{Sylow-$p$-Untergruppe} von $G$. Wir schreiben $n_p(G)$ für die Anzahl der
Sylow-$p$-Untergruppen.
\end{definition}

\begin{theorem}[Erster Sylow-Satz]
\label{thm:sylow1}
Sei $G$ eine endliche Gruppe, $p$ eine Primzahl und $p^k \mid |G|$. Dann besitzt $G$
eine Untergruppe der Ordnung $p^k$. Insbesondere existieren Sylow-$p$-Untergruppen.
\end{theorem}

\begin{proof}[Beweisidee]
Durch Induktion über $|G|$. Falls $p \nmid [G:Z(G)]$, dann gilt $p \mid |Z(G)|$, und
der Satz von Cauchy liefert ein Element der Ordnung $p$ im Zentrum. Im allgemeinen Fall
verwendet man die Klassengleichung
$|G| = |Z(G)| + \sum [G:C_G(x)]$ und wendet Induktion auf eine echte Untergruppe an.
\end{proof}

\begin{theorem}[Zweiter Sylow-Satz]
\label{thm:sylow2}
Je zwei Sylow-$p$-Untergruppen von $G$ sind konjugiert: Für beliebige
Sylow-$p$-Untergruppen $P$ und $Q$ gibt es ein $g \in G$ mit $Q = gPg^{-1}$.
\end{theorem}

\begin{theorem}[Dritter Sylow-Satz]
\label{thm:sylow3}
Die Anzahl $n_p = n_p(G)$ der Sylow-$p$-Untergruppen erfüllt:
\begin{enumerate}[label=(\roman*)]
  \item $n_p \equiv 1 \pmod{p}$,
  \item $n_p \mid [G : P] = |G|/p^k$.
\end{enumerate}
Insbesondere gilt $n_p = 1$ genau dann, wenn die eindeutige Sylow-$p$-Untergruppe ein
Normalteiler in $G$ ist.
\end{theorem}

\begin{example}[Gruppen der Ordnung 15]
Sei $|G| = 15 = 3 \cdot 5$. Dann gilt $n_3 \mid 5$ und $n_3 \equiv 1 \pmod{3}$, also
$n_3 = 1$. Analog gilt $n_5 = 1$. Beide Sylow-Untergruppen sind Normalteiler, daher
$G \cong \mathbb{Z}_3 \times \mathbb{Z}_5 \cong \mathbb{Z}_{15}$. Es gibt genau eine
Gruppe der Ordnung 15, und sie ist zyklisch.
\end{example}

\begin{corollary}[Satz von Cauchy]
Ist $p$ eine Primzahl und $p \mid |G|$, so enthält $G$ ein Element der Ordnung $p$.
\end{corollary}

%% ============================================================
\section{Auflösbare und nilpotente Gruppen}
\label{sec:aufloesbar}
%% ============================================================

\subsection{Kompositionsreihen und Jordan--Hölder}

\begin{definition}[Kompositionsreihe]
Eine \emph{Kompositionsreihe} für $G$ ist eine Kette von Untergruppen
\[
  1 = G_0 \trianglelefteq G_1 \trianglelefteq \cdots \trianglelefteq G_n = G,
\]
sodass jeder Quotient $G_{i+1}/G_i$ einfach ist (nichttrivial und ohne echte Normalteiler).
Die Quotienten $G_{i+1}/G_i$ heißen \emph{Kompositionsfaktoren}.
\end{definition}

\begin{theorem}[Satz von Jordan--Hölder]
\label{thm:jordan_hoelder}
Besitzt $G$ eine Kompositionsreihe, so haben je zwei Kompositionsreihen von $G$ dieselbe
Länge, und ihre Kompositionsfaktoren stimmen bis auf Permutation und Isomorphismus überein.
\end{theorem}

\subsection{Auflösbare Gruppen}

\begin{definition}[Ableitung und Auflösbarkeit]
Die \emph{Ableitung} (oder \emph{Kommutatorenreihe}) von $G$ ist
\[
  G = G^{(0)} \geq G^{(1)} \geq G^{(2)} \geq \cdots,
  \quad G^{(i+1)} = [G^{(i)}, G^{(i)}],
\]
wobei $[H, K]$ die von allen Kommutatoren $[h,k] = hkh^{-1}k^{-1}$ erzeugte Untergruppe
bezeichnet. Die Gruppe $G$ heißt \emph{auflösbar}, wenn $G^{(n)} = 1$ für ein $n$ gilt.
\end{definition}

\begin{theorem}
Eine endliche Gruppe $G$ ist genau dann auflösbar, wenn alle Kompositionsfaktoren in einer
beliebigen Kompositionsreihe von $G$ zyklisch von Primzahlordnung sind.
\end{theorem}

\begin{theorem}[Satz von Abel--Ruffini, 1824]
\label{thm:abel_ruffini}
Für $n \geq 5$ ist die symmetrische Gruppe $S_n$ \emph{nicht} auflösbar (da $A_n$ einfach
und nicht-abelsch ist). Folglich kann eine generische Polynomgleichung vom Grad $\geq 5$
nicht durch Radikale gelöst werden.
\end{theorem}

\begin{proof}[Zusammenhang mit der Galois-Theorie]
Ein Polynom $f \in \mathbb{Q}[x]$ ist genau dann durch Radikale lösbar, wenn seine
Galois-Gruppe $\mathrm{Gal}(f)$ auflösbar ist. Für das allgemeine Polynom vom Grad
$n \geq 5$ gilt $\mathrm{Gal}(f) \cong S_n$, das die einfache, nicht-abelsche Untergruppe
$A_n$ enthält. Daher ist $S_n$ für $n \geq 5$ nicht auflösbar.
\end{proof}

\subsection{Nilpotente Gruppen}

\begin{definition}[Untere Zentralreihe und Nilpotenz]
Die \emph{untere Zentralreihe} von $G$ ist
\[
  G = \gamma_1(G) \geq \gamma_2(G) \geq \cdots,
  \quad \gamma_{i+1}(G) = [G, \gamma_i(G)].
\]
$G$ heißt \emph{nilpotent} von \emph{Klasse} $c$, wenn $\gamma_{c+1}(G) = 1$, aber
$\gamma_c(G) \neq 1$.
\end{definition}

\begin{theorem}
Jede nilpotente Gruppe ist auflösbar. Eine endliche Gruppe ist genau dann nilpotent,
wenn sie das direkte Produkt ihrer Sylow-$p$-Untergruppen ist. Insbesondere ist jede
endliche $p$-Gruppe (Gruppe von Primzahlpotenzordnung) nilpotent.
\end{theorem}

%% ============================================================
\section{Einfache Gruppen}
\label{sec:einfach}
%% ============================================================

\begin{definition}[Einfache Gruppe]
Eine Gruppe $G$ mit $|G| > 1$ heißt \emph{einfach}, wenn ihre einzigen Normalteiler $1$
und $G$ selbst sind.
\end{definition}

\begin{theorem}[$A_n$ ist einfach für $n \geq 5$]
\label{thm:an_einfach}
Die alternierende Gruppe $A_n$ ist für alle $n \geq 5$ einfach.
\end{theorem}

\begin{proof}[Beweisidee]
Jeder nichttriviale Normalteiler von $A_n$ muss einen $3$-Zyklus enthalten
(unter Verwendung der Tatsache, dass $A_n$ von $3$-Zyklen erzeugt wird, und mithilfe
von Konjugationsklassen-Argumenten). Jeder Normalteiler, der einen $3$-Zyklus enthält,
enthält alle $3$-Zyklen und stimmt daher mit $A_n$ überein.
\end{proof}

\subsection{Projektiv-spezielle lineare Gruppen}

\begin{definition}
Für eine Primzahlpotenz $q = p^r$ und $n \geq 1$ ist die \emph{projektiv-spezielle
lineare Gruppe} $\mathrm{PSL}(n, q)$ der Quotient
$\mathrm{SL}(n, q) / Z(\mathrm{SL}(n,q))$, wobei $\mathrm{SL}(n,q)$ die Gruppe der
$n \times n$-Matrizen über $\mathbb{F}_q$ mit Determinante $1$ ist.
\end{definition}

\begin{theorem}
Für $n \geq 2$ und $q$ eine Primzahlpotenz ist $\mathrm{PSL}(n, q)$ einfach, außer für
$\mathrm{PSL}(2, 2) \cong S_3$ und $\mathrm{PSL}(2, 3) \cong A_4$.
Die Gruppen $\mathrm{PSL}(2, q)$ für Primzahlpotenzen $q \geq 4$ bilden unendliche
Familien nicht-abelscher einfacher Gruppen.
\end{theorem}

\subsection{Satz von Feit--Thompson und CFSG}

\begin{theorem}[Satz von Feit--Thompson, 1963]
\label{thm:feit_thompson}
Jede endliche Gruppe ungerader Ordnung ist auflösbar.
\end{theorem}

\begin{remark}
Der Satz von Feit--Thompson füllt einen ganzen Band des \emph{Pacific Journal of
Mathematics} (255 Seiten). Sein Beweis ist ein Meisterwerk der Darstellungstheorie,
Charaktertheorie und kombinatorischen Gruppentheorie. Als unmittelbare Konsequenz
hat jede nicht-abelsche endliche einfache Gruppe gerade Ordnung.
\end{remark}

\begin{theorem}[Klassifikation endlicher einfacher Gruppen (CFSG), 2004]
\label{thm:cfsg}
Jede endliche einfache Gruppe gehört zu einer der folgenden Familien:
\begin{enumerate}[label=(\roman*)]
  \item \textbf{Zyklische Gruppen} $\mathbb{Z}_p$ für Primzahlen $p$.
  \item \textbf{Alternierende Gruppen} $A_n$ für $n \geq 5$.
  \item \textbf{Gruppen vom Lie-Typ}: 16 unendliche Familien, einschließlich
        $\mathrm{PSL}(n,q)$, $\mathrm{PSp}(2n,q)$, $\mathrm{P}\Omega^{\pm}(n,q)$,
        $G_2(q)$, $F_4(q)$, $E_6(q)$, $E_7(q)$, $E_8(q)$ und verdrehte Varianten
        ${}^2A_n(q)$, ${}^2D_n(q)$, ${}^3D_4(q)$, ${}^2G_2(q)$, ${}^2F_4(q)$,
        ${}^2E_6(q)$, ${}^2B_2(q)$.
  \item \textbf{26 sporadische Gruppen}: darunter die Monstergruppe $\mathbb{M}$
        (Ordnung $\approx 8 \times 10^{53}$), die Babymonster und die 20 Mitglieder
        der ``Glücklichen Familie'', die mit der Monstergruppe durch Moonshine
        verwandt sind.
\end{enumerate}
\end{theorem}

\begin{remark}
Der Beweis der CFSG wurde um 1983 mit wesentlichen Beiträgen von Gorenstein, Lyons,
Solomon, Aschbacher, Smith und vielen anderen abgeschlossen. Ein überarbeiteter und
vollständiger Beweis, die ``Klassifikation zweiter Generation'', wurde 2004 in einer
12-bändigen Reihe fertiggestellt. Der Beweis umfasst über 10.000 Seiten in Hunderten
von Zeitschriftenartikeln.
\end{remark}

\begin{theorem}[Monstergruppe]
Die \emph{Monstergruppe} $\mathbb{M}$ ist die größte sporadische einfache Gruppe.
Ihre Ordnung beträgt:
\begin{align*}
  |\mathbb{M}| &= 2^{46} \cdot 3^{20} \cdot 5^9 \cdot 7^6 \cdot 11^2 \cdot 13^3
                  \cdot 17 \cdot 19 \cdot 23 \cdot 29 \cdot 31 \cdot 41 \cdot 47
                  \cdot 59 \cdot 71\\
               &\approx 8{,}08 \times 10^{53}.
\end{align*}
\end{theorem}

%% ============================================================
\section{Anwendungen}
\label{sec:anwendungen}
%% ============================================================

\subsection{Galois-Theorie}

Gruppentheorie und Körpertheorie sind durch die Galois-Theorie verknüpft, eines der
schönsten Kapitel der Mathematik.

\begin{theorem}[Hauptsatz der Galois-Theorie]
\label{thm:galois_hauptsatz}
Sei $L/K$ eine endliche galoissche Erweiterung mit Galois-Gruppe $\mathrm{Gal}(L/K)$.
Es gibt eine inklusionsumkehrende Bijektion zwischen Zwischenkörpern
$K \subseteq F \subseteq L$ und Untergruppen $H \leq \mathrm{Gal}(L/K)$:
\[
  F \;\longleftrightarrow\; H = \mathrm{Gal}(L/F), \quad
  [L : F] = |H|, \quad [F : K] = [\mathrm{Gal}(L/K) : H].
\]
Dabei ist $F/K$ genau dann galoisch, wenn $H \trianglelefteq \mathrm{Gal}(L/K)$;
in diesem Fall gilt $\mathrm{Gal}(F/K) \cong \mathrm{Gal}(L/K)/H$.
\end{theorem}

\begin{example}
Der Zerfällungskörper von $x^4 - 2$ über $\mathbb{Q}$ ist $L = \mathbb{Q}(\sqrt[4]{2},i)$
mit $\mathrm{Gal}(L/\mathbb{Q}) \cong D_4$ (Diedergruppe der Ordnung 8). Der
Untergruppenverband von $D_4$ entspricht bijektiv den 10 Zwischenkörpern zwischen
$\mathbb{Q}$ und $L$.
\end{example}

\subsection{Kristallographische Gruppen}

\begin{definition}
Eine \emph{Raumgruppe} ist eine diskrete Untergruppe $\Gamma$ der euklidischen Gruppe
$\mathrm{E}(3)$, sodass der Quotient $\mathbb{R}^3/\Gamma$ kompakt ist. Die
\emph{Punktgruppe} ist das Bild von $\Gamma$ in $\mathrm{O}(3)$.
\end{definition}

\begin{theorem}[Kristallographische Einschränkung]
In Dimension 2 kann eine kristallographische Gruppe (Wandgruppe) nur Drehsymmetrien der
Ordnung $2, 3, 4$ oder $6$ enthalten. Es gibt genau $17$ Wandgruppen und genau
$230$ Raumgruppen in drei Dimensionen.
\end{theorem}

\begin{remark}
Die 230 Raumgruppen, unabhängig von Fedorov (1891) und Schönflies (1891) klassifiziert,
bilden die mathematische Grundlage der Kristallographie. Moderne Kristallstrukturbestimmung
mittels Röntgenbeugung nutzt diese Klassifikation auf jeder Stufe.
\end{remark}

\subsection{Quantenmechanik: \texorpdfstring{$\mathrm{SU}(2)$}{SU(2)} und
            \texorpdfstring{$\mathrm{SO}(3)$}{SO(3)}}

\begin{definition}
Die \emph{spezielle unitäre Gruppe}
$\mathrm{SU}(2) = \{A \in \mathrm{M}_{2}(\mathbb{C}) : A^\dagger A = I,\, \det A = 1\}$
und die \emph{Drehgruppe}
$\mathrm{SO}(3) = \{R \in \mathrm{M}_3(\mathbb{R}) : R^T R = I,\, \det R = 1\}$
sind Lie-Gruppen, die durch die doppelte Überlagerung $\mathrm{SU}(2) \twoheadrightarrow
\mathrm{SO}(3)$ mit Kern $\{\pm I\}$ verknüpft sind.
\end{definition}

\begin{theorem}
Die irreduziblen komplexen Darstellungen von $\mathrm{SU}(2)$ sind durch nicht-negative
Halbganzzahlen $j \in \{0, \frac{1}{2}, 1, \frac{3}{2}, 2, \ldots\}$ parametrisiert
und haben Dimension $2j+1$. Die Darstellungen steigen genau dann zu $\mathrm{SO}(3)$
ab, wenn $j$ eine ganze Zahl ist.
\end{theorem}

\begin{remark}
In der Quantenmechanik ist $j$ die \emph{Spin}-Quantenzahl. Elektronen, Protonen und
Neutronen sind Spin-$\frac{1}{2}$-Teilchen und transformieren unter der 2-dimensionalen
Darstellung von $\mathrm{SU}(2)$. Photonen sind Spin-$1$-Teilchen und transformieren
unter $\mathrm{SO}(3)$. Die Wigner-$3j$-Symbole beim Drehimpulskoppeln sind die
Clebsch--Gordan-Koeffizienten der Zerlegung
$j_1 \otimes j_2 = |j_1 - j_2| \oplus \cdots \oplus (j_1 + j_2)$.
\end{remark}

\subsection{Kryptographie: Gruppen des diskreten Logarithmus}

\begin{definition}
Sei $G = \langle g \rangle$ eine endliche zyklische Gruppe der Ordnung $n$.
Das \emph{diskrete Logarithmusproblem (DLP)} lautet: Gegeben $h \in G$, finde
$x \in \mathbb{Z}$ mit $g^x = h$.
\end{definition}

\begin{remark}
Für geeignet gewählte Gruppen (z.\,B.\ Untergruppen von Primzahlordnung in
$(\mathbb{Z}/p\mathbb{Z})^*$ oder Punkte auf elliptischen Kurven über endlichen Körpern)
ist kein polynomieller klassischer Algorithmus für das DLP bekannt. Diese Schwierigkeit
bildet die Sicherheitsgrundlage für Protokolle wie Diffie--Hellman-Schlüsselaustausch,
ElGamal-Verschlüsselung und ECDSA-Signaturen. Shors Quantenalgorithmus löst das DLP in
polynomieller Zeit auf einem Quantencomputer, was den Übergang zur Post-Quanten-Kryptographie
motiviert.
\end{remark}

%% ============================================================
\section{Schluss}
\label{sec:schluss}
%% ============================================================

Wir haben den Kern der endlichen Gruppentheorie durchlaufen: von den grundlegenden Axiomen
über die Isomorphiesätze und die Klassifikation endlicher abelscher Gruppen bis hin zu den
mächtigen Werkzeugen der Sylow-Theorie und den Strukturresultaten über auflösbare und
nilpotente Gruppen. Der Satz von Jordan--Hölder zeigt, dass Kompositionsfaktoren ein
gruppentheoretisches Invariante sind, und der Satz von Abel--Ruffini zeigt, wie die
Auflösbarkeit von Gruppen direkt in die Lösbarkeit von Gleichungen durch Radikale
übersetzt.

Die Klassifikation endlicher einfacher Gruppen (CFSG), eine der größten kollektiven
Leistungen in der Geschichte der Mathematik, liefert eine vollständige Liste der
``Atome'', aus denen alle endlichen Gruppen aufgebaut sind. Der Satz von Feit--Thompson,
dass jede endliche Gruppe ungerader Ordnung auflösbar ist, ist eine tiefe Voraussetzung
für die CFSG.

Anwendungen der Gruppentheorie durchziehen Mathematik und Physik. Die Galois-Theorie
nutzt Gruppentheorie, um antike Fragen über Polynomgleichungen zu beantworten.
Kristallographische Gruppen regeln die Struktur aller bekannten Kristalle. Die
Darstellungstheorie von $\mathrm{SU}(2)$ und $\mathrm{SO}(3)$ bildet die Grundlage der
Quantentheorie des Drehimpulses und des Spins. Zyklische Gruppen und ihre Analoga bilden
das Rückgrat der modernen Public-Key-Kryptographie.

Offene Richtungen umfassen die Monstrous-Moonshine-Vermutungen, die die Monstergruppe mit
Modulformen verbinden (Borcherds' Beweis der Conway--Norton-Vermutung, 1992), die
Baum--Connes-Vermutung, die gruppen-$K$-Theorie mit Analysis verknüpft, und das
Langlands-Programm, das gruppen-, zahlen- und geometrie-theoretische Strukturen durch
die Darstellungstheorie reduktiver Gruppen vereint.

%% ============================================================
\section*{Literatur}
%% ============================================================

\begin{thebibliography}{99}

\bibitem{Artin2011}
M.~Artin, \emph{Algebra}, 2.~Aufl., Pearson, 2011.

\bibitem{Dummit2004}
D.~S.~Dummit und R.~M.~Foote, \emph{Abstract Algebra}, 3.~Aufl.,
Wiley, Hoboken, NJ, 2004.

\bibitem{Rotman2012}
J.~J.~Rotman, \emph{Advanced Modern Algebra}, 3.~Aufl., American Mathematical
Society, Providence, RI, 2010.

\bibitem{FeitThompson1963}
W.~Feit und J.~G.~Thompson, ``Solvability of groups of odd order'',
\emph{Pacific J.\ Math.}\ \textbf{13} (1963), Nr.~3, 775--1029.

\bibitem{Sylow1872}
L.~Sylow, ``Théorèmes sur les groupes de substitutions'',
\emph{Math.\ Ann.}\ \textbf{5} (1872), 584--594.

\bibitem{Gorenstein1982}
D.~Gorenstein, \emph{Finite Simple Groups: An Introduction to Their
Classification}, Plenum Press, New York, 1982.

\bibitem{Lang2002}
S.~Lang, \emph{Algebra}, 3.~Aufl., Springer, New York, 2002.

\bibitem{Hall1959}
M.~Hall, \emph{The Theory of Groups}, Macmillan, New York, 1959.

\bibitem{Serre1977}
J.-P.~Serre, \emph{Linear Representations of Finite Groups}, Springer,
New York, 1977.

\bibitem{Aschbacher2004}
M.~Aschbacher, R.~Lyons, S.~D.~Smith und R.~Solomon,
\emph{The Classification of Finite Simple Groups: Groups of Characteristic 2 Type},
Amer.\ Math.\ Soc., 2011.

\end{thebibliography}

\end{document}
